\begin{theorem}[Explicit real involutions and all fixed fibres]\hypertarget{tag-02PF}{}\label{mapsI:src:A.3.2}
\StatementRecord{mapsI:src:A.3.2}{proved}{02PF}{\SeriesRef{fam:src:B.4.2}, \SeriesRef{mapsI:real-eq-4}, \SeriesRef{mapsI:real-eq-2}, \SeriesRef{fam:src:A.3.1}, \SeriesRef{mapsI:real-eq-1}}

On the ordinary period charts the formulas
\begin{equation}
 \sigma_c[z,\zeta]=[\iota z,\bar\zeta+b_c(\iota z)],
 \qquad \sigma_c^- =\rho(1/2)\sigma_c
 \label{mapsI:real-eq-3}
\end{equation}
extend to antiholomorphic involutions of the compact threefold.
With these representatives their complete fixed-set data are as follows.
Here $2T^2$ means the disjoint union of two real two-tori; every displayed
circle in a double curve avoids its two triple-point endpoints.
\begin{center}\small
\begin{tabular}{cclccc}
\toprule
$w$&involution&image in $\RP^1$&on $S_1$&on $S_2$&on $W$\\
\midrule
$-ir$, $r>0$&$\sigma_c$&all $\RP^1$&$2T^2$&$2T^2$&$T^2\sqcup S^1_{D_2}\sqcup S^1_{D_{12}}$\\
$-ir$, $r>0$&$\sigma_c^-$&$[1,\infty]$&$\varnothing$&$T^2$&$S^1_{D_1}\sqcup S^1_{D_2}$\\
$ir$, $r>0$&$\sigma_c$&$\varnothing$&$\varnothing$&$\varnothing$&$\varnothing$\\
$ir$, $r>0$&$\sigma_c^-$&$[\infty,1]$ through $0$&$2T^2$&$T^2$&$S^1_{D_1}\sqcup S^1_{D_{12}}$\\
\bottomrule
\end{tabular}
\end{center}
The first interval is the closure of $t>1$; the second is the closure
of $t<1$. On every open real arc contained in the indicated image the
real regular fibre is $2T^2$. No triple point is real for these structures.
\end{theorem}



\begin{proof}
\begin{enumerate}[label=\arabic*.,leftmargin=*]
\item Start with the reflection $(z,\zeta)\mapsto(\iota z,-\bar\zeta)$
and the integral matrix $H$ of \eqref{I:eq:real-reflection-matrix}.
Proposition~\ref{fam:src:B.4.2} proves its period compatibility and gives
the target affine datum \eqref{I:eq:real-reflected-affine-data}.
The normalization from this
target at parameter $c_0=-\bar c$ to the canonical datum at $c_0+1$
has linear part $-I$ and translation
\begin{equation}
 b^{\rm norm}_{c_0}=\tfrac13\Pi_{c_0}g+(2L_F-4L_G)e.
 \hypertarget{tag-02P5}{}\label{mapsI:real-eq-4}
\end{equation}
Its period defects at the two finite generators are
$(0,-2,2,0)^t$ and $(0,0,2,0)^t$. To check them, the logarithmic increments
are $-2/3$ and $1$, and one multiplies the two displayed matrices $A_j$
by $g/3$. The finite difference vectors in the complex translation space are $2e$ and $-4e$, so subtracting
their logarithmic collars leaves holomorphic translations on both full
root discs. At the cusp
$2L_F-4L_G=s/3-(2\pi i)^{-1}\log(1-q)$ and
$\Pi g=(6\mu,-s+C_{21})^t$. The logarithmic terms cancel and give the
holomorphic toric unit
$(2\mu,C_{21}/3-(2\pi i)^{-1}\log(1-q))$.
Thus \eqref{mapsI:real-eq-4} extends over all three fillings, with its inverse.

\item For $c=1/2+iy$ compose this normalization with the initial reflection
and then with the constant translation $e/6$. This gives exactly the
first line of \eqref{mapsI:real-eq-2}. For $c=-1/2+iy$, additionally close the parameter
shift by the following explicit even-shift formula with $h=-2$,
\[
 b_{h,c_0}=-\frac h{12}\Pi_{c_0}g+
 \left(h(L_F-L_G)-\frac{h^2}{24}\right)e.
\]
A final constant translation $-e/4$ gives the second line of \eqref{mapsI:real-eq-2}.
The displayed translation represents the even shift of Proposition~\ref{fam:src:A.3.1}. Its finite-chart corrections are holomorphic. For its cusp expression, put $\ell=L_F-L_G=-s/12+\kappa_0(q)$ and
$C_{c,21}(q)=c+C_{0,21}(q)$, and set
$D_0(q)=C_{0,21}(q)-12\kappa_0(q)$.
The explicit even-shift translation becomes
\[
 b_{h,c}=
 \left(-h\mu/2,-h(c+D_0)/12-h^2/24\right)^t.
\]
The logarithmic terms have cancelled exactly. Exponentiating both
coordinates gives holomorphic nonzero torus multipliers on the whole cusp
disc; multiplication by these units extends over every toric stratum.
Compatibility with the period deck group is the integral identity
$\exp(2\pi i h\lambda_1)=1$ for $h\in2\mathbb Z$ and
$\lambda_1\in\mathbb Z$. The finite-chart defects are checked by the
same displayed $A_j$ and root logarithms used for the normalization.

\item Over $t>1$ we have
$\Re\mu=1/2$, $\Re\tau=\Re\beta^{\rm part}=0$ and the logarithms in
\eqref{mapsI:real-eq-1} are purely imaginary. The two squares are respectively the
integral period translations $2\widehat w$ and $3\widehat w$.
The square is a holomorphic automorphism and is consequently the
identity everywhere. This proves the involution claim.

\item We next give exact equations for all real fibres. Put $a=\Re c$ and
$R=-H+2aE$. In the chosen branch the reflection identifications on
$t>1$, $0<t<1$, $t<0$ are respectively $1,g_2,g_0^{-1}$.
They follow by reflecting the chosen upper-half-plane paths; equivalently
the root arguments are $(0,0)$, $(0,\pi/4)$ and $(\pi/3,\pi/4)$.
The three fibre involution matrices are
\begin{equation}
 R_>=R,\qquad R_m=A_2^{-1}R,\qquad R_<=M_0R.
 \hypertarget{tag-02P6}{}\label{mapsI:real-eq-5}
\end{equation}
In marked real coordinates let $r_c(t)=\Pi_c(t)^{-1}b_c(t)$,
where $\Pi_c^{-1}$ is the real inverse. The fixed tori are exactly
\begin{equation}
 \{x\in\R^4/\Lambda:(I-R_\bullet)x=
                       r_c(t)+\epsilon d/2\pmod\Lambda\},
 \qquad\epsilon=0,1.
 \label{mapsI:real-eq-6}
\end{equation}
This specifies the cosets, including their parameter dependence. Their
identity subgroups are $T_{\langle w,d\rangle}$ on both outer arcs and
$T_{\langle u+w,d\rangle}$ on the middle arc. When nonempty there are
two cosets. Relative to any one solution their offsets are
$(g+u)/2$ on the right arc and $g/2$ on the other two arcs.

\item Existence of fixed points is decided by integral affine congruences. An affine involution
$x\mapsto R_\bullet x+\beta$ has a fixed point if and only if its
square translation $(I+R_\bullet)\beta$ belongs to
$(I+R_\bullet)\Lambda$. In \eqref{mapsI:real-eq-2} the square translation is
\[
 (I+R_\bullet)\alpha g+2\Re S\,d.
\]
For $c=1/2+iy$ the values $2\Re S$ on the three arcs are
$-1/3,-4/3,-2/3$; for $c=-1/2+iy$ they are $1/2,0,0$.
The resulting integral square vectors are respectively
$(2w,0,2w)$ and $(3w,d,3w)$. Adding $d/2$ to the involution adds $d$
to each square vector. Taking column Hermite normal forms of
$I+R_\bullet$ gives exactly the nonempty entries in the table.
All three $I-R_\bullet$ have nonzero Smith factors $1,2$.
This proves the coset counts and offsets. The Hermite and Smith
computations are listed in \cref{tab:I:programs}. The left-arc lift
$g_0^{-1}$ follows from the reflected upper-half-plane path above.

\item At $t=0$ the root degree is odd. Conjugation by the degree-three deck
group acts transitively on the three reflection lifts; the centralizer
of a reflection in that deck group is trivial.
Thus the quotient has exactly the same two real tori as the adjacent
nonempty real regular fibres. In the empty cases there are none. This
also follows by solving the limiting affine equations on the full
multiplicity-three root quotient.

\item At $t=1$ write the affine deck map as
$G(x)=A_2x+v_2/4$. The real root involution has matrix $R_>$ and
translation
\[
 b_-=(1/3,0,0,-1/6)^t,
 \qquad b_+=(1/2,0,0,1/4)^t
\]
for the two parameter signs. The real locus of the quotient is the
quotient of the fixed sets of $G^kh$, $0\leq k<4$. Its even and odd
lift parities are identified internally by $G$, while $G^2$ preserves
each parity. Direct integer congruences show that the nonempty lift
parities are: both for negative imaginary $w$ and $\epsilon=0$; only even for
negative imaginary $w$ and $\epsilon=1$; neither for positive imaginary $w$ and $\epsilon=0$;
only odd for positive imaginary $w$ and $\epsilon=1$. Every nonempty lift has
two tori, exchanged by $G^2$. Each parity therefore contributes one
torus downstairs. The congruence enumeration uses denominator $24$:
translation denominators divide $12$, and the nonzero Smith factors are
$1,2$, so every component has such a representative. The component
permutation is included in the exact enumeration listed in
\cref{tab:I:programs}. The normal
real line has fourth power of sign positive for the even parity and
negative for the odd parity. This establishes the two one-sided
endpoints at $t=1$ recorded in the table.

\item We compute the cusp fixed set. Put
$A=C_{12}(0)$; the cusp normalization gives
\[
 \Re\mu(0)=1/2,\quad A\in i\R,\quad C_{21}(0)=c-A,
 \quad C_{11}(0)=6\mu(0),\quad C_{22}(0)=\mu(0).
\]
In a toric lift the involution is
$x\mapsto U/\bar x$, with
\begin{equation}
\begin{array}{c|cc}
 &b_1(0)&b_2(0)\\\hline
 w=-ir&2\mu(0)&C_{21}(0)/3-1/6\\
 w=ir&3\mu(0)&C_{21}(0)/2+1/4
\end{array}
\qquad U_j=e^{2\pi i b_j(0)}.
\hypertarget{tag-02P7}{}\label{mapsI:real-eq-7}
\end{equation}
Thus $U_2>0$ in both rows, whereas $U_1>0$ in the first and $U_1<0$
in the second. Replacing $\sigma_c$ by $\sigma_c^-$ negates $U_2$.
The smooth central orbit has a real torus precisely when both units
are positive, with equations $|x_j|^2=U_j$.

\item For an edge $v$, set $n=(-v_2,v_1)^t$. The period translate bringing
the inverted edge back is $B_0n=v$. Its curve coordinate $z=x^n$ has
involution
\begin{equation}
 z\longmapsto K_v/\bar z,
 \qquad K_v=U^n e^{2\pi i n^tC(0)n}.
 \hypertarget{tag-02P8}{}\label{mapsI:real-eq-8}
\end{equation}
The signs for $(D_1,D_2,D_{12})$ are $(-,+,+)$ for negative $w$
and $(-,-,-)$ for positive $w$. Negating $U_2$ reverses the first and
third signs, but not the second. A positive sign gives the single
circle $|z|^2=K_v$; a negative sign gives no fixed point. Central
inversion exchanges the two triangle classes of the periodic fan, so
it exchanges the two triple points. These calculations prove every
cusp entry. Finally a nonempty fixed three-manifold cannot be confined
to the special fibres: a holomorphic function constant on an open
piece of a totally real submanifold is locally constant. This also
proves that the positive-$w$ representative with no regular real
fibres is entirely fixed-point free.
\end{enumerate}
\end{proof}
